\section{Tổng quan}
\label{sec:tongquan}

\subsection{Giới thiệu}

Griductive Solver là đồ án thứ hai của học phần CSC14003 -- Nhập môn
Trí tuệ nhân tạo. Đồ án xây dựng một trò chơi suy luận trên lưới và một
tác tử có khả năng giải tự động dựa trên logic mệnh đề. Trên bàn cờ
$N\times N$, mỗi ô là một thẻ nhân vật có một trong hai trạng thái
\texttt{CRIMINAL} (tội phạm) hoặc \texttt{INNOCENT} (vô tội). Mục tiêu
của người chơi là xác định trạng thái của mọi nhân vật và mở toàn bộ thẻ
bằng các kết luận suy ra từ thông tin công khai.

Khi bắt đầu một màn, một số thẻ đã mở và công bố cả trạng thái của nhân
vật lẫn manh mối do nhân vật đó sở hữu. Một manh mối có thể xác định trực
tiếp trạng thái của một ô, liên hệ trạng thái giữa hai ô hoặc giới hạn
số tội phạm trong một vùng. Chẳng hạn, nếu hàng đầu tiên có đúng một
tội phạm và đã biết B1 là tội phạm, người chơi có thể kết luận các ô
còn lại trên hàng đó đều vô tội. Sau khi một kết luận được chấp nhận,
trò chơi mở thẻ tương ứng. Manh mối mới xuất hiện có thể giúp chứng minh
trạng thái của những ô trước đó chưa xác định được.

Đặc điểm của bài toán là thông tin được bổ sung qua từng nước đi. Một
kết luận về ô chưa mở chỉ hợp lệ khi trạng thái đó đúng trong mọi cách
gán còn phù hợp với dữ kiện công khai tại thời điểm xét. Nếu cả hai
trạng thái của ô vẫn có thể xảy ra, người chơi cần tìm kết luận ở ô
khác. Nếu điều này xảy ra với mọi ô chưa mở, tác tử báo \texttt{STUCK}.
Vì vậy, hai câu hỏi cần được kiểm tra riêng là toàn bộ manh mối có xác
định một nghiệm duy nhất hay không và từ các thẻ mở ban đầu có thể suy
luận từng bước đến khi hoàn tất hay không.

\subsection{Mục tiêu đồ án}

\begin{enumerate}
    \item \textbf{Mô hình hóa trò chơi bằng logic mệnh đề.} Biểu diễn
    trạng thái mỗi nhân vật bằng một biến Boolean, xây dựng ngữ nghĩa
    cho các loại manh mối và tập hợp các dữ kiện đã công bố thành cơ sở
    tri thức tại mỗi bước.
    \item \textbf{Chuyển cơ sở tri thức sang CNF và kiểm tra SAT.} Xây
    dựng bộ mã hóa để chuyển manh mối thành dạng chuẩn hội
    (\emph{conjunctive normal form}, CNF). Cài đặt thuật toán DPLL để
    kiểm tra tính khả thỏa của công thức bằng lan truyền đơn vị, phân
    nhánh và quay lui.
    \item \textbf{Xây dựng tác tử suy luận.} Dùng DPLL để kiểm tra xem
    trạng thái đối lập với kết luận có khả thỏa cùng cơ sở tri thức hay
    không. Từ đó, tác tử phân loại từng ô, đề xuất nước đi có thể chứng
    minh và nhận diện trường hợp thiếu dữ kiện hoặc cơ sở tri thức mâu
    thuẫn. Tác tử nhận trạng thái công khai từ Game Engine, còn đáp án
    đầy đủ và các manh mối chưa mở thuộc dữ liệu nội bộ của Engine.
    \item \textbf{Tích hợp suy luận vào trò chơi.} Dùng chung quy tắc
    xác nhận để xử lý kết luận do người chơi gửi và nước đi do tác tử đề
    xuất. Cung cấp giao diện để chọn màn, thao tác trên bàn, xem gợi ý
    và theo dõi quá trình giải tự động.
    \item \textbf{Kiểm chứng và đánh giá.} Kiểm tra tính duy nhất của
    nghiệm, mô phỏng quá trình mở thẻ từ trạng thái ban đầu và đo chi
    phí suy luận trên các màn có kích thước và mức độ khó khác nhau.
\end{enumerate}

\subsection{Phạm vi và đóng góp}

Chương trình được viết bằng Python với giao diện Tkinter. Bộ dữ liệu
gồm chín màn lưu dưới dạng JSON, với ba màn cho mỗi kích thước
$3\times3$, $4\times4$ và $5\times5$. Hai màn mẫu $3\times3$ và
$4\times4$ minh họa chuỗi suy luận từ một thẻ mở ban đầu. Các màn còn
lại cung cấp những tổ hợp manh mối khác nhau để kiểm tra khả năng
suy luận trên nhiều trạng thái bàn. Giao diện hỗ trợ chơi thủ công,
chọn màn, nạp tệp JSON, khởi động lại, lấy gợi ý, giải tự động và dừng
quá trình giải tự động.

Bộ mã hóa hỗ trợ sáu loại manh mối \texttt{FACT}, \texttt{SAME},
\texttt{DIFFERENT}, \texttt{EXACTLY}, \texttt{AT\_LEAST} và
\texttt{AT\_MOST}. Các manh mối đếm có thể áp dụng trên hàng, cột, tập
ô lân cận hoặc danh sách ô chỉ định. Hai loại vùng mở rộng cho phép
xét bốn góc của bàn và phần giao giữa tập ô lân cận của hai nhân vật.
Nhóm tự cài đặt bộ mã hóa CNF, thuật toán DPLL và tác tử suy luận.
Thiết kế tách việc xử lý bàn cờ trong Game Engine, điều phối nước đi
trong Game Controller và suy luận logic trong Logic Agent.

Ngoài luồng chơi, chương trình có bộ kiểm định màn để kiểm tra cấu trúc,
xác định tính duy nhất của nghiệm khi xét toàn bộ manh mối và kiểm tra
khả năng giải tuần tự khi manh mối được công bố theo luật chơi. Bộ công
cụ thực nghiệm ghi lại số biến, số mệnh đề CNF, số bước suy luận, số lần
gọi SAT, các thống kê của DPLL và thời gian giải. Kết quả được trình bày
theo từng màn và theo kích thước bàn, trong phạm vi chín màn của đồ án.

Các mục tiếp theo trình bày mô hình bài toán, quy tắc mã hóa CNF,
thuật toán DPLL, cách Logic Agent chọn nước đi, kiến trúc giao diện và
kết quả thực nghiệm.
